\section{Discussion}
\label{sec:discussion}

\subsection{Why seed sensitivity is low}

The low seed sensitivity (CV $< 2\%$ for 48 of 50 states) is attributable
to the structure of census tract adjacency graphs.  These graphs are planar,
sparse, and have bounded diameter relative to the grid.  The planar separator
theorem~\citep{lipton1979separator,b6paper} implies that optimal separators
are approximately equally distributed around the graph's ``waist,'' so many
seeds find similar or equivalent separators.

This contrasts with the results of DeFord et al.~\citep{deford2019recombination}
on unrestricted redistricting plans, where seed sensitivity (measured by
partisan variance) can be much higher.  The difference is that METIS
optimises a specific objective (minimum edge cut) rather than sampling
uniformly from the feasible set; the objective concentrates solutions
near a small number of near-optimal configurations.

\subsection{The outlier states: Georgia and North Carolina}

Both outlier states ($k = 14$) exhibit higher variance for a structural
reason: 14 is neither a power of two nor a highly composite number.
The bisection tree for $k = 14$ requires splits of 7+7, then 4+3 and 4+3,
then 2+2, 2+1, 2+2, 2+1 --- four levels with asymmetric splits at every
level above the bottom.  Each asymmetric split introduces variance because
the optimal separator for a 7:3 split differs from a 1:1 split.

The geographic complexity of Georgia and North Carolina compounds this
effect: both states have multiple ``pinch points'' --- geographic features
(rivers, ridges) that constrain separator placement --- creating multiple
near-optimal configurations that different seeds explore.

\subsection{Implications for the DIA's single-seed specification}

The results support the DIA's single-seed specification for 48 of 50 states.
For the two high-variance states (GA, NC), the ConvergenceSweep procedure
(B.16) provides additional assurance: running seeds until $T = 600$
non-improving seeds certifies convergence to a local minimum within
the low-variance regime of the solution space.

A key legal point: the DIA seed is derived from the census release
identifier via SHA-256, making it \emph{unknowable before the census
is published}.  No party can precompute which seed produces a favorable
outcome.  Combined with the low seed sensitivity result, this makes the
DIA seed specification robust to adversarial manipulation.

\subsection{Limitations}

The partisan analysis uses 2020 presidential election data, which may
not reflect future electoral patterns.  The 10,000-seed sweep is large
but not exhaustive: the true minimum-edge-cut seed may not be among
the 10,000 tested.  However, the $\varepsilon$-ball analysis
($83\%$ of seeds within 5\% of the observed minimum) suggests that
10,000 seeds is sufficient to characterise the distribution.

The results apply to the B.2 edge-weighted configuration.  Unweighted
configurations or alternative weight functions (county-sticky, B.10) may
show different seed sensitivity profiles; this is an open empirical question.
